% Generated draft: requires mathematical and source review.
\documentclass[11pt]{article}
\usepackage[margin=25mm]{geometry}
\usepackage{fontspec}
\setmainfont{Latin Modern Roman}
\usepackage{amsmath,amssymb,mathtools,unicode-math}
\setmathfont{Latin Modern Math}
\usepackage{enumitem,needspace,xurl,hyperref}
\setlength{\parindent}{0pt}
\setlength{\parskip}{5pt}
\setlength{\emergencystretch}{3em}
\newcommand{\N}{\mathbb{N}}
\newcommand{\Z}{\mathbb{Z}}
\newcommand{\Q}{\mathbb{Q}}
\newcommand{\R}{\mathbb{R}}
\newcommand{\C}{\mathbb{C}}
\newcommand{\F}{\mathbb{F}}
\newcommand{\A}{\mathbb{A}}
\newcommand{\PP}{\mathbb{P}}
\newcommand{\E}{\mathbb{E}}
\newcommand{\eps}{\varepsilon}
\newcommand{\dd}{\,\mathrm{d}}
\begin{document}
\section*{UnsolvedMath ID: 30003382}
\textbf{Source problem code:} catalogue entry
\subsection{Definitions and mathematical statement}
The supplied record describes a research programme rather than a single fully specified conjecture. Let \(G\) denote a weighted graph, let \(X_G\) denote its space of admissible states, and write \(f_G\) for the PDE-inspired graph functional called \(f\) in the source. Let \(X\) denote a continuum space of admissible states and \(F\) its proposed continuum counterpart. Here PDE means partial differential equation. The graph class, weight conventions, state spaces, functional formulas, continuum domain, constraints, and boundary conditions are not supplied; these symbols are placeholders for data that must be specified in each instance.

The requested programme is to identify such data and establish correspondences addressing the following questions:
\begin{enumerate}
\item Which structural properties of \(F\) admit precisely formulated analogues for \(f_G\), and under which hypotheses do those analogues hold?
\item For which sequences of weighted graphs \((G_n)_{n\geq 1}\), choices of weights and scaling, and identifications of graph states with continuum states is there a convergence result of the schematic form
\[
 f_{G_n}\xrightarrow{\mathcal C}F \qquad (n\to\infty)?
\]
Here \(\mathcal C\) denotes a notion of convergence to be specified, including the treatment of the varying spaces \(X_{G_n}\). The display does not assert convergence or select a graph-limit regime, topology, or variational convergence notion.
\item How can algorithms for minimizing \(f_G\), with their admissible sets and constraints specified, be used in data analysis and image processing, and what mathematical conclusions support those uses?
\item Which additional network problems admit PDE-inspired methods, and which additional PDE or variational systems admit meaningful and useful graph analogues?
\end{enumerate}
The source does not specify the structural properties to be covered, the meaning of usefulness, or quantifiers over a defined class of functionals. Consequently, it does not support a universal assertion for every pair \((f_G,F)\). A precise theorem or conjecture requires an explicit instance or class together with the missing hypotheses and convergence definitions.
\subsection{Short English statement}
The problem asks for graph counterparts of continuum variational properties, conditions for graph functionals to converge to continuum functionals, applications of graph minimization to data and images, and further PDE-inspired network models. It is an open-ended programme: the supplied record does not specify enough mathematical data for a single quantified conjecture. Its supplied status is partially solved, subject to the dated source qualifications below.
\subsection{Sources}
\begin{itemize}
\item Local UnsolvedMath registry: supplied record with permanent numeric ID \(30003382\), exact source problem code \texttt{catalogue entry}, title ``PDE-Inspired Graph Functionals and Their Continuum Limits,'' in the Oberwolfach Reports Open Problems collection. No individual proposer is supplied; the proposed year is recorded as 2017.
\item Supplied report citation: \emph{Emerging Developments in Interfaces and Free Boundaries} (2018), Oberwolfach Reports, pp.\ 267--338, DOI \texttt{10.4171/owr/2017/6}. \url{https://doi.org/10.4171/owr/2017/6}. The record describes Oberwolfach Reports as documenting mathematical workshops and their open-problem discussions; no specific problem page or passage is supplied.
\item Supplied literature reference, retaining its supplied attribution and title: ``Bertozzi et al., Graph Ginzburg--Landau Gamma-convergence.'' \url{https://arxiv.org/abs/1204.5220}.
\item Supplied literature reference: ``Convergence of the Graph Allen--Cahn Scheme.'' \url{https://link.springer.com/article/10.1007/s10955-017-1772-4}.
\item Literature assessment dated exactly \texttt{2026-08-21}: the supplied curation labels the problem partially solved and reports established graph analogues and continuum limits for total variation, Ginzburg--Landau, Allen--Cahn, and related models, while reporting no known general correspondence covering all listed properties and PDE-inspired network problems. This is supplied dated literature triage derived from the OWR curation, not an independently verified assessment. Time-sensitive status and the cited claims require independent verification.
\item The registry labels the statement \texttt{corrected\_verified} and says that mathematical content was preserved while notation, typography, or omitted definitions were repaired. This is a supplied editorial qualification, not independent mathematical verification of this draft.
\end{itemize}
\end{document}
